% part: intuitionistic-logic
% chapter: introduction
% section: bhk-interpretation

\documentclass[../../../include/open-logic-section]{subfiles}

\begin{document}

\olfileid{int}{int}{bhk}

\olsection{ব্রাউয়ার--হাইটিং--কলমোগোরভ ব্যাখ্যা}

\begin{editorial}
  BHK ব্যাখ্যা দিয়ে স্বজ্ঞাবাদী বচনের বৈধতা প্রমাণগুলি বিভ্রান্তিকর;
  এগুলি আরও ভালোভাবে ব্যাখ্যা করা দরকার।
\end{editorial}

স্বজ্ঞাবাদী সংযোজকগুলির একটি অনানুষ্ঠানিক নির্মাণমূলক ব্যাখ্যা আছে,
যাকে সাধারণত ব্রাউয়ার--হাইটিং--কলমোগোরভ ব্যাখ্যা বলা হয়। এতে
``নির্মাণ'' ধারণাটি ব্যবহৃত হয়; একে নির্মাণমূলক প্রমাণ হিসেবে ভাবতে
পারো। (আনুষ্ঠানিক !!{derivation} ব্যবস্থায় !!a{derivation}-এর ধারণার
সঙ্গে গুলিয়ে না ফেলতে BHK ব্যাখ্যায় আমরা ``প্রমাণ'' শব্দটি ব্যবহার
করি না।) এই স্বজ্ঞাত ধারণার ভিত্তিতে BHK ব্যাখ্যা স্বজ্ঞাবাদী
সংযোজকগুলির অর্থ বোঝায়।

\begin{enumerate}
\item ধরে নিই, একটি পারমাণবিক উক্তির নির্মাণ বলতে কী বোঝায় তা আমরা জানি।
\item $!A_1 \land !A_2$-এর নির্মাণ হলো একটি জোড়া $\tuple{M_1, M_2}$,
  যেখানে $M_1$ হলো~$!A_1$-এর নির্মাণ এবং $M_2$ হলো~$!A_2$-এর নির্মাণ।
\item $!A_1 \lor !A_2$-এর নির্মাণ হলো একটি জোড়া $\tuple{s, M}$,
  যেখানে হয় $s$ হলো~$1$ এবং $M$ হলো~$!A_1$-এর নির্মাণ, অথবা $s$ হলো~$2$
  এবং $M$ হলো~$!A_2$-এর নির্মাণ।
\item $!A \lif !B$-এর নির্মাণ হলো এমন একটি অপেক্ষক, যা~$!A$-এর
  একটি নির্মাণকে~$!B$-এর একটি নির্মাণে রূপান্তর করে।
\item $\lfalse$-এর (অসংগতির) কোনো নির্মাণ নেই।
\item $\lnot !A$-কে $!A \lif \lfalse$-এর সমার্থক বলে সংজ্ঞায়িত
  করা হয়। অর্থাৎ, $\lnot !A$-এর নির্মাণ হলো এমন একটি অপেক্ষক, যা~$!A$-এর
  নির্মাণকে~$\lfalse$-এর নির্মাণে রূপান্তর করে।
\end{enumerate}

\begin{ex}
উদাহরণ হিসেবে $\lnot \lfalse$ নাও। এর নির্মাণ হলো এমন একটি অপেক্ষক,
যা ইনপুট হিসেবে $\lfalse$-এর যেকোনো নির্মাণ পেলে আউটপুট হিসেবে
$\lfalse$-এর একটি নির্মাণ দেয়। স্পষ্টতই অভেদ অপেক্ষক~$\Id{}$ এমন একটি
নির্মাণ: $\lfalse$-এর কোনো নির্মাণ~$M$ দেওয়া হলে,
$\Id{}(M) = M$-ও $\lfalse$-এর একটি নির্মাণ দেয়।
\end{ex}

সাধারণভাবে $\lnot !A$-এর অর্থ ``$!A$-এর নির্মাণ অসম্ভব''।

\begin{ex}
যেকোনো বচন~$!A$-এর জন্য $!A \lif \lnot \lnot !A$ প্রমাণ করা যাক;
এটি হলো $!A \lif ((!A \lif \lfalse) \lif \lfalse)$। নির্মাণটি হতে
হবে এমন একটি অপেক্ষক~$f$, যা $!A$-এর একটি নির্মাণ~$M$ পেলে
$(!A \lif \lfalse) \lif \lfalse$-এর একটি নির্মাণ~$f(M)$ ফেরায়।
$f$ কীভাবে $(!A \lif \lfalse) \lif \lfalse$-এর নির্মাণ গড়ে তা
দেখি: এমন একটি অপেক্ষক $g$ সংজ্ঞায়িত করতে হবে, যা ইনপুট হিসেবে
$!A \lif \lfalse$-এর একটি নির্মাণ~$h$ পেলে আউটপুট হিসেবে
$\lfalse$-এর একটি নির্মাণ দেয়। $g$-কে এভাবে সংজ্ঞায়িত করা যায়:
ইনপুট~$h$-কে আগে পাওয়া $!A$-এর নির্মাণ~$M$-এ প্রয়োগ করো। যেহেতু
$h$-এর আউটপুট~$h(M)$ হলো $\lfalse$-এর নির্মাণ, $M$ যদি~$!A$-এর
নির্মাণ হয়, তবে $f(M)(h) = h(M)$ হলো~$\lfalse$-এর নির্মাণ।
\end{ex}

\begin{ex}
$\lnot(!A \land \lnot !A)$-এর, অর্থাৎ $(!A
\land (!A \lif \lfalse)) \lif \lfalse$-এর একটি নির্মাণ দিই। এটি
এমন একটি অপেক্ষক~$f$, যা ইনপুট হিসেবে $!A \land (!A \lif \lfalse)$-এর
একটি নির্মাণ~$M$ পেলে $\lfalse$-এর একটি নির্মাণ দেয়। $!B_1 \land
!B_2$ সংযোজনের নির্মাণ হলো একটি জোড়া $\tuple{N_1, N_2}$, যেখানে
$N_1$ হলো~$!B_1$-এর এবং $N_2$ হলো~$!B_2$-এর নির্মাণ। $p_1$ ও $p_2$
অপেক্ষক এমনভাবে সংজ্ঞায়িত করা যায়, যাতে তারা $!B_1 \land !B_2$-এর
নির্মাণ থেকে যথাক্রমে $!B_1$ ও $!B_2$-এর নির্মাণ উদ্ধার করে:
\begin{align*}
  p_1(\tuple{N_1, N_2}) &= N_1\\
  p_2(\tuple{N_1, N_2}) &= N_2
\end{align*}
$f$-এর কাজটি হলো: প্রথমে সে নিজের ইনপুট~$M$-এ $p_1$ প্রয়োগ করে।
তাতে~$!A$-এর একটি নির্মাণ মেলে। তারপর~$M$-এ $p_2$ প্রয়োগ করে,
যাতে $!A \lif \lfalse$-এর একটি নির্মাণ মেলে। এই নির্মাণটি আবার
একটি অপেক্ষক~$p_2(M)$: ইনপুট হিসেবে~$!A$-এর নির্মাণ পেলে সে
$\lfalse$-এর নির্মাণ দেয়। অন্য কথায়, $p_2(M)$-কে $p_1(M)$-এ
প্রয়োগ করলে $\lfalse$-এর একটি নির্মাণ পাই। সুতরাং
$f(M) = p_2(M)(p_1(M))$ সংজ্ঞায়িত করা যায়।
\end{ex}

\begin{ex}
$((!A \land !B) \lif !C) \lif (!A \lif
(!B \lif !C))$-এর একটি নির্মাণ দিই; অর্থাৎ, এমন একটি অপেক্ষক $f$
দিই, যা $(!A \land !B) \lif !C$-এর একটি নির্মাণ~$g$-কে
$(!A \lif (!B \lif !C))$-এর নির্মাণে রূপান্তর করে।
% BN-SRC-394: the codomain of g is !C, matching the displayed implication.
$g$ নির্মাণটি নিজেই একটি অপেক্ষক ($!A \land !B$-এর নির্মাণ থেকে
$!C$-এর নির্মাণে)। আর আউটপুট $f(g)$ হলো এমন একটি অপেক্ষক~$h_g$,
যা $!A$-এর নির্মাণ থেকে অপেক্ষকে যায়; সেই অপেক্ষকগুলি~$!B$-এর
নির্মাণকে~$!C$-এর নির্মাণে পাঠায়।

ঠিক আছে, ব্যাপারটি জটিল। আমাদের এমন একটি অপেক্ষক $h_g$ নির্মাণ
করতে হবে, যা ইনপুট~$g$-এর জন্য $f$-এর আউটপুট হবে। $h_g$-এর ইনপুট
হলো~$!A$-এর একটি নির্মাণ~$M$। $h_g(M)$-এর আউটপুট হতে হবে এমন একটি
অপেক্ষক~$k_M$, যা~$!B$-এর নির্মাণ~$N$ থেকে~$!C$-এর নির্মাণে যায়।
ধরি, $k_{g,M}(N) = g(\tuple{M, N})$। মনে রেখো, $\tuple{M,
  N}$ হলো~$!A \land !B$-এর একটি নির্মাণ। তাই $k_{g,M}$ হলো
$!B \lif !C$-এর নির্মাণ: এটি~$!B$-এর নির্মাণ~$N$-কে~$!C$-এর
নির্মাণে পাঠায়। এখন ধরি, $h_g(M) = k_{g,M}$। এই অপেক্ষক~$!A$-এর
নির্মাণ~$M$-কে $!B \lif !C$-এর নির্মাণ~$k_{g,M}$-এ পাঠায়। এবার
ধরি, $f(g) = h_g$। এই অপেক্ষক $(!A \land !B) \lif !C$-এর
নির্মাণ~$g$-কে $!A \lif (!B \lif !C)$-এর নির্মাণে পাঠায়। উফ!{}
\end{ex}

$!A \lor \lnot !A$ উক্তিটি বর্জিত-মধ্যম নীতি নামে পরিচিত। কোনো
বিশেষ $!A$-এর জন্য (যেমন, $\lfalse \lor
\lnot \lfalse$) আমরা এটি প্রমাণ করতে পারি, কিন্তু সাধারণভাবে
পারি না। কারণ স্বজ্ঞাবাদী বিকল্পের ক্ষেত্রে দুটি বিকল্পাংশের
কোনো একটির নির্মাণ চাই; অথচ এমন উক্তি আছে যা এখন পর্যন্ত প্রমাণও
করা যায়নি, খণ্ডনও করা যায়নি (যেমন, গোল্ডবাখের অনুমান)। তবে
বর্জিত-মধ্যম নীতিটিকেও খণ্ডন করা যায় না: অর্থাৎ, $\lnot \lnot (!A
\lor \lnot !A)$ সত্য।

\begin{ex}
$\lnot \lnot (!A \lor \lnot !A)$ প্রমাণ করতে এমন একটি অপেক্ষক~$f$
চাই, যা $\lnot(!A \lor \lnot !A)$-এর, অর্থাৎ $(!A
\lor (!A \lif \lfalse)) \lif \lfalse$-এর একটি নির্মাণকে
$\lfalse$-এর নির্মাণে রূপান্তর করে। অন্য কথায়, এমন $f$ চাই যাতে
$g$ যদি $\lnot(!A \lor \lnot !A)$-এর নির্মাণ হয়, তবে $f(g)$ হয়
$\lfalse$-এর নির্মাণ।

ধরো, $g$ হলো $\lnot(!A \lor \lnot !A)$-এর একটি নির্মাণ; অর্থাৎ,
এমন একটি অপেক্ষক, যা $!A \lor \lnot !A$-এর নির্মাণকে $\lfalse$-এর
নির্মাণে পাঠায়। $!A \lor \lnot !A$-এর নির্মাণ হলো একটি জোড়া
$\tuple{s, M}$: হয় $s = 1$ এবং $M$ হলো $!A$-এর নির্মাণ, অথবা
$s = 2$ এবং $M$ হলো $\lnot !A$-এর নির্মাণ। ধরি, $h_1$ এমন একটি
অপেক্ষক, যা $!A$-এর নির্মাণ~$M_1$-কে $!A \lor \lnot !A$-এর নির্মাণে
পাঠায়:
% BN-SRC-393: h_1 receives M_1, so the first injection uses M_1 rather than M_2.
এটি $M_1$-কে $\tuple{1, M_1}$-এ পাঠায়। আর ধরি, $h_2$ এমন একটি
অপেক্ষক, যা $\lnot !A$-এর নির্মাণ~$M_2$-কে $!A \lor \lnot !A$-এর
নির্মাণে পাঠায়: এটি $M_2$-কে $\tuple{2, M_2}$-এ পাঠায়।

ধরি, $k$ হলো $\comp{h_1}{g}$: ইনপুট হিসেবে~$!A$-এর একটি নির্মাণ
পেলে এটি $\lfalse$-এর একটি নির্মাণ ফেরায়; অর্থাৎ, এটি
$!A \lif \lfalse$ বা $\lnot !A$-এর একটি নির্মাণ। এবার ধরি,
$l$ হলো $\comp{h_2}{g}$। এটি এমন একটি অপেক্ষক, যা $\lnot !A$-এর
নির্মাণ পেলে $\lfalse$-এর নির্মাণ দেয়। যেহেতু $k$ হলো
$\lnot !A$-এর একটি নির্মাণ, তাই $l(k)$ হলো $\lfalse$-এর একটি নির্মাণ।

সব মিলিয়ে আমরা দেখিয়েছি, কীভাবে $\lnot(!A \lor \lnot !A)$-এর
একটি নির্মাণ~$g$ থেকে $\lfalse$-এর একটি নির্মাণ পাওয়া যায়। অর্থাৎ,
যে অপেক্ষক $f$, $\lnot(!A \lor \lnot !A)$-এর নির্মাণ~$g$-কে
$\lfalse$-এর নির্মাণ~$l(k)$-এ পাঠায়, সেটি
$\lnot\lnot(!A \lor \lnot !A)$-এর একটি নির্মাণ।
\end{ex}

দেখতেই পাচ্ছ, BHK ব্যাখ্যা দিয়ে !!{formula}s-এর স্বজ্ঞাবাদী
বৈধতা দেখানো দ্রুতই দুরূহ ও বিভ্রান্তিকর হয়ে ওঠে। সৌভাগ্যবশত,
স্বজ্ঞাবাদী যুক্তিবিদ্যার জন্য আরও ভালো !!{derivation} ব্যবস্থা
এবং আরও নির্ভুল অর্থবিষয়ক ব্যাখ্যা আছে।
\end{document}
